% right_panel.tex — FIG_PLAN_2026-09-19 §3.3: assemble the four modules into
% the right panel. Only positioning + panel borders + title bars + the two
% cross-module relations happen here; module internals are untouched.
%
% Native module CONTENT canvases (read off each modX_*.tex \path rectangle;
% this is the area each module's own coordinates were designed against):
%   A 2.94 x 3.03   B 6.53 x 3.03   C 4.735 x 3.03   D 4.735 x 3.03
% Row content width: A+0.25+B = 9.72 = C+0.25+D  (matches, no stretching
% needed).
%
% v3 (user: titles must span the full width of each panel, centred, and
% nothing may sit under them). Rather than shrink/reflow each module's own
% (already tightly-packed, already-gated) internal content to make room for
% a full-width bar, this ADDS a new strip on top of each panel exclusively
% for the bar and leaves every module's own coordinate system untouched --
% each module still only knows about its own 3.03-tall content area; the
% bar lives in the extra \TitleH above it, entirely at this assembly layer.
% v3: 0.36 -> 0.46. Needed so panel A's two-line title (see \PanelTitle
% above) has room; applied to all 4 panels uniformly so every title bar
% stays the same height (the row layout only knows one \TitleH).
\def\TitleH{0.46}          % title-bar strip height, same for all 4 panels
\def\ModH{3.03}            % each module's own content height (unchanged)
\pgfmathsetmacro{\PanelH}{\ModH+\TitleH}   % = 3.39, full panel box height

\newcommand{\PanelBox}[3]{% #1=(x,y) origin (box bottom-left)  #2=width  #3=height
  \path[panel] #1 rectangle ++(#2,#3);
}
% NOTE: deliberately uses a shifted scope + plain-brace pgfmath, not calc's
% ($...$) syntax -- calc's $ delimiters read as math-mode $ to
% fig_symbol_audit.py (G2), which then flagged \TitleH as an unmatched
% notation symbol. This local geometry macro is not paper notation, so it
% should never appear in a G2 report; this form keeps it out of one.
\newcommand{\PanelTitle}[4]{% #1=(x,y) bottom-left of the title strip
                             % #2=width  #3=fill colour  #4=centred text
\begin{scope}[shift={#1}]
  \path[rounded corners=3pt,fill=#3] (0,0) rectangle (#2,\TitleH);
  \node[font=\bfseries\scriptsize,text=white,align=center] at ({0.5*#2},{0.5*\TitleH}) {#4};
\end{scope}}
% v3: the plan's new title text ("A Finite-memory state" etc.) is longer
% than the old text this macro was tuned for, and panel A's own width
% (2.94cm, unchanged -- the 2x2 layout is on the "immutable" list) is the
% narrowest of the four; at one line, \scriptsize, it clipped against the
% panel's own rounded corner in the render. G3 sets \scriptsize as a FLOOR
% ("minimum font size"), so the fix is a two-line title for A alone (see
% \TitleH below), not a smaller font for anyone. align=center (added here,
% affects all four uniformly, so still one shared macro, not a per-panel
% override) is what makes A's explicit "\\" below actually wrap.

% Placement (each module's own local (0,0) origin, in this panel's frame;
% row 1 sits \PanelH above row 2 plus the usual 0.25cm gap):
%   A -> (0.00, \PanelH+0.25)      B -> (3.19, \PanelH+0.25)
%   C -> (0.00, 0.00)              D -> (4.985, 0.00)
\pgfmathsetmacro{\RowOneY}{\PanelH+0.25}   % = 3.64

% ---- panel borders (full height, including the title strip), drawn first
% so the title fill and module content both sit on top of them ----
\PanelBox{(0.00,\RowOneY)}{2.94}{\PanelH}
\PanelBox{(3.19,\RowOneY)}{6.53}{\PanelH}
\PanelBox{(0.00,0.00)}{4.735}{\PanelH}
\PanelBox{(4.985,0.00)}{4.735}{\PanelH}

% ---- title bars, in the strip ABOVE each module's own content area (at
% local y=\ModH within that module's frame) -- a NEW band, not carved out
% of space any module's internals already use, so nothing needed to move ----
% v3 (2026-09-19, post CP2-redo): titles renamed to match the plan's
% §2.2 wording, which in turn matches §3.1-3.3's own subsection titles.
\PanelTitle{(0.00,\RowOneY+\ModH)}{2.94}{clModA}{A Finite-memory\\state}
\PanelTitle{(3.19,\RowOneY+\ModH)}{6.53}{clModB}{B Koopman dynamics model}
\PanelTitle{(0.00,\ModH)}{4.735}{clModC}{C Learning from data}
\PanelTitle{(4.985,\ModH)}{4.735}{clModD}{D Predictive selection}

% ---- module content, each in its own shifted scope (origin unchanged --
% only the panel box and title strip grew, not the module's own frame) ----
\begin{scope}[shift={(0.00,\RowOneY)}]  \path (0,0) rectangle (2.94,3.03);

% ---- 1. the window slides out of a stream of observations ----
\MemRow{(1.30,2.68)}{7}{0.20}
\DimBraceT{(1.90,2.68)}{(2.70,2.68)}{$\symMemDim$}
\node[mathlab,anchor=north,inner sep=1pt] at (2.60,2.47) {$\symAttr$};
\draw[thinflow] (2.30,2.48) -- (1.80,2.09);

% ---- 2. block algebra: the window's first entry IS the current value ----
\RowSelector{(0.65,1.76)}{4}{0.22}
\node[mathlab,inner sep=1pt] at (1.61,1.65) {$\cdot$};
\MemColV{(1.69,2.09)}{4}{0.22}
\node[mathlab,inner sep=1pt] at (2.02,1.65) {$=$};
\ScalarCell{(2.13,1.76)}{tok1}{$\symAttr$}
\node[mathlab,anchor=north,inner sep=1pt] at (1.09,1.54) {$\symSel$};
\node[mathlab,anchor=north,inner sep=1pt] at (1.80,1.21) {$\symMem$};

% ---- 3. why one observation is not enough ----
\MiniAxes{(0.10,0.32)}{1.85}{0.65}
\draw[draw=bad,line width=0.6pt] (0.10,0.93) -- (0.72,0.80) -- (1.30,0.70) -- (1.95,0.645);
\draw[draw=ours,line width=0.6pt] (0.10,0.36) -- (0.72,0.48) -- (1.30,0.58) -- (1.95,0.645);
\draw[-{Stealth[length=1.6mm]},dash pattern=on 1.4pt off 1.2pt,line width=0.5pt,draw=bad] (1.95,0.645) -- (2.18,0.50);
\draw[-{Stealth[length=1.6mm]},dash pattern=on 1.4pt off 1.2pt,line width=0.5pt,draw=ours] (1.95,0.645) -- (2.18,0.79);
\fill[black] (1.95,0.645) circle (0.025);
\node[minicap,inner sep=1pt] at (1.47,0.16) {same $\symAttr$};
   \end{scope}
\begin{scope}[shift={(3.19,\RowOneY)}]  % modB_koopman.tex — module B: the command-conditioned Koopman dynamics model.
% v3 (2026-09-19, post CP2-redo): eq:koopman_model is stated directly on the
% reading window s_t (no encoder/decoder in the main text), so the main chain
% now runs left to right entirely in s-space:
%   s_t -> [K, (L+1)x(L+1)] -> (+) -> shat_{t+1} -> [C selector] -> yhat_{t+1}
% with the command column B (+ the scalar c_t) entering the sum from below --
% this branch is still the point of the figure -- and the bias b entering
% from above. The one thing kept from the lifting story (3.1 P6) is drawn as
% a single faded, optional detour arcing over the top of the main chain,
% through a de-emphasised encoder/decoder pair; the paper reports the
% identity lifting (the solid chain below), so the arc+nets are opacity 0.45
% and never the thing a reader's eye lands on first. The two labels that
% belong to the arc (E_phi, and the "optional learned lifting" caption) are
% kept at full opacity so they stay legible (G4) -- only the drawn shapes
% (line + nets) are faded, per the plan's "opacity=0.45" on the arc itself.
% Canvas 6.53 x 3.03; top-left strip is the title bar (added in right_panel).
\path (0,0) rectangle (6.53,3.03);

% ---------------- 1. finite-memory state s_t --------------------------
% top-left y=2.09, same local coordinate module A uses for its own s_t
% column bottom -- the cross-module arrow in right_panel.tex depends on
% this and is NOT changed here. L+1 is already labelled once, in A; per
% the plan's "no dimension appears twice" rule this module does not repeat
% \DimBraceR for it.
\MemColV{(0.70,2.09)}{4}{0.22}
\node[mathlab,anchor=north] at (0.81,1.19) {$\symMem$};
\draw[thinflow] (0.92,1.75) -- (1.45,1.75);

% ---------------- 2. the transition matrix K ---------------------------
\MatBlock{(1.45,2.19)}{4}{4}{0.22}{clLatentD}
\node[mathlab,anchor=north] at (1.89,1.29) {$\symK$};
\DimBraceT{(1.45,2.21)}{(2.33,2.21)}{$\symMemDim\times{}\symMemDim$}
% the "{}" above is a no-op in the typeset output (same as v20's render);
% it exists only so the G2 audit script's naive string-substitution macro
% expander does not read "\times" immediately followed by \symMemDim's
% own expansion "L+1" as the single (nonexistent) control word "\timesL".
% Real TeX tokenizes \times and \symMemDim as separate control words
% regardless, so this changes nothing LaTeX actually renders -- it is the
% same class of audit-script false positive the plan's risk table already
% flags for \symResid's \lVert, fixed the same way it prescribes: work
% around it in the fragment, do not touch the script, note it here.
\draw[thinflow] (2.33,1.75) -- (2.86,1.75);

% ---------------- 3. summing node, bias (above) and command (below) ----
\node[mathlab,circle,draw=black!60,line width=0.5pt,fill=white,
      inner sep=0pt,minimum size=3.2mm] (Bsum) at (3.02,1.75) {$+$};
% bias branch, from above -- kept small (cell 0.10) and labelled to its
% WEST (not above) so the column's own top can stay well below the lifting
% arc that passes over it (see step 6): with this column, its top sits at
% y=2.39, and the arc/nets are kept at y>=2.58, a 0.19cm clearance, comfortably
% above the plan's >=0.15cm minimum.
\MatBlock{(2.97,2.39)}{4}{1}{0.10}{black}
\node[mathlab,anchor=east] at (2.94,2.19) {$\symBias$};
\draw[thinflow] (3.02,1.99) -- (3.02,1.91);
% command branch, from below: still the point of the figure, so it keeps
% the full-size cell (0.22) and the accent colour. Column is 4x1 (matches
% B's real shape, (L+1)x1). Label placed to the column's WEST (like b's,
% above) rather than stacked below it -- the first attempt stacked $B$'s
% label directly above the "command" minicap at the old (3-row) offsets,
% and re-using those same y's for a taller 4-row column left only a hair
% of clearance, so the two merged into unreadable "comBand" text.
\MatBlock{(2.91,1.25)}{4}{1}{0.22}{ours}
\node[mathlab,anchor=east] at (2.88,0.81) {$\symBcol$};
\ScalarCell{(3.15,1.03)}{ours}{$\symCmd$}
\draw[oursflow] (3.02,1.25) -- (3.02,1.59);
\node[minicap,anchor=south,text=ours] at (3.02,0.15) {command};

% ---------------- 4. predicted state shat_{t+1} -------------------------
\draw[thinflow] (3.18,1.75) -- (3.70,1.75);
\MemColV{(3.70,2.09)}{4}{0.22}
\node[mathlab,anchor=west] at (3.94,2.05) {$\symMemHat$};

% ---------------- 5. readout: selection row C, then the scalar yhat ----
\draw[thinflow] (3.92,1.75) -- (4.45,1.75);
\RowSelector{(4.45,1.86)}{4}{0.22}
\node[mathlab,anchor=north] at (4.89,1.62) {$\symSel$};
\draw[thinflow] (5.33,1.75) -- (5.86,1.75);
\ScalarCell{(5.86,1.86)}{tok1}{$\symYhat$}

% ---------------- 6. optional learned lifting: one faded detour --------
% 3.1 P6: "the same affine step can be written on learned features E_phi(s_t)
% instead of on the window; this paper reports the identity lifting." Drawn
% once, as a complete arc from s_t's own column top back to shat_{t+1}'s own
% column top, through a de-emphasised encoder/decoder pair -- never a bare
% "id" bypass (that macro belonged to the retired xi-space chain and is not
% used here).
%   The nets sit in the two GAPS either side of the main chain's busy
% middle (Enc over the s_t->K arrow, Dec just before shat_{t+1}), not
% directly above K or the bias column -- the first attempt put Enc right
% above K and its label collided with K's own (L+1)x(L+1) brace label,
% because both are centred on K's own x-span; here Enc's centre (x~1.18)
% is well clear of the brace (x 1.45-2.33). Only the thin horizontal
% connecting line, not a net's full height, has to clear the bias column
% (top y=2.39) and K's brace (top y~2.36) in between -- at y=2.70 it clears
% both by >=0.3cm, comfortably past the plan's 0.15cm minimum.
\begin{scope}[opacity=0.45]
  \draw[bypass] (0.81,2.09) to[out=70,in=200] (1.04,2.80);
  \begin{scope}[shift={(1.10,2.80)},scale=0.24,transform shape]
    \EncNet{mbLiftE}{(0,0)}
  \end{scope}
  \draw[bypass] (1.378,2.80) -- (3.314,2.80);
  \begin{scope}[shift={(3.372,2.80)},scale=0.24,transform shape]
    \DecNet{mbLiftD}{(0,0)}
  \end{scope}
  \draw[bypass] (3.65,2.80) to[out=-20,in=100] (3.81,2.09);
\end{scope}
% E_phi appears exactly once in the whole figure: here, under the one
% de-emphasised encoder that draws it (G10). The decoder that follows it is
% deliberately left unlabelled -- 3.1 P6 introduces one new symbol (E_phi),
% not a paired decoder symbol, and the main text never names one.
%   Two earlier revisions tried to clear the K-dimension label
% ($(L+1)\times(L+1)$, centred over K) by moving THIS label up and by
% shrinking/raising the nets -- both still rendered as merged run-on text
% (only visible in the PNG, not from the coordinates: the brace macro's
% decoration bumps its text noticeably higher than the raw brace-line y
% suggests). Fixed instead by getting this label off that shared row
% entirely: it now sits to the WEST of the encoder, vertically centred on
% the net (y=2.80), far above the brace's own row (y<=2.45).
\node[mathlab,anchor=east] at (0.98,2.80) {$\symEnc$};
\node[minicap,anchor=south west] at (4.55,1.05) {optional learned lifting};
  \end{scope}
\begin{scope}[shift={(0.00,0.00)}]      % modC_learning.tex — module C: learning from interaction data (§3.2).
% v6 (balance/layout pass, 2026-09-19): canvas narrowed to 3.40 so both
% rows split at the same x; content stacked in three tiers, top to bottom,
% all connectors orthogonal:
%   1. readings recorded from the frozen model, with the sliding window,
%   2. the transition tuples the window cuts out,
%   3. the least-squares fit (eq:identification): predicted successor
%      against the recorded one, one residual link, fitted parameters
%      drawn as the SAME glyph shapes as the unknowns (4x4 block, 4x1
%      columns) -- v4 drew Khat as a single cell, a scalar glyph.
% p_theta sits at the data source: it produced the readings and takes no
% part in the fit (v4 had it in a corner, half outside the panel).
\path (0,0) rectangle (3.40,3.03);

% ---- 1. readings from the frozen model; sliding window -------------------
\SnowIcon[0.80]{(0.30,2.72)}
\node[mathlab,anchor=north,inner sep=1pt] at (0.30,2.54) {$\symLLM$};
\draw[thinflow] (0.47,2.72) -- (0.58,2.72);
\MemRow{(0.60,2.80)}{8}{0.16}
\WinBracket{(0.60,2.80)}{4}{0.16}{ours}
\WinBracket{(0.76,2.80)}{4}{0.16}{black!50}
\node[mathlab,anchor=west] at (1.91,2.72) {$\symAttr$};
\draw[thinflow] (0.66,2.62) -- (0.66,2.24);

% ---- 2. transition tuples (three cards, back to front) -------------------
\foreach \dx in {0.12,0.06,0.00}{%
  \begin{scope}[shift={(\dx,\dx)}]
    \MemColV{(0.62,2.10)}{4}{0.10}
    \ScalarCell{(0.75,1.96)}{ours}{}
    \MemColV{(1.00,2.10)}{4}{0.10}
  \end{scope}}
\node[minicap,anchor=west] at (1.32,2.30) {sliding window};
\node[mathlab,anchor=west] at (1.32,1.95) {$\symTrans$};
\draw[thinflow] (0.88,1.70) -- (0.88,1.36);

% ---- 3. least squares: predicted successor vs recorded, baseline 1.10 ---
\MatBlock{(0.30,1.34)}{4}{4}{0.12}{clLatentD}
\node[mathlab,anchor=north] at (0.54,0.82) {$\symK$};
\MemColV{(0.82,1.34)}{4}{0.12}
\node[mathlab,anchor=north] at (0.88,0.82) {$\symMem$};
\node[mathlab] at (1.04,1.10) {$+$};
\MatBlock{(1.12,1.34)}{4}{1}{0.12}{ours}
\node[mathlab,anchor=north] at (1.18,0.82) {$\symBcol$};
\ScalarCellT{(1.27,1.21)}{ours}{$\symCmd$}
\node[mathlab] at (1.58,1.10) {$+$};
\MatBlock{(1.66,1.34)}{4}{1}{0.12}{black}
\node[mathlab,anchor=north] at (1.72,0.82) {$\symBias$};
\draw[lossline] (1.84,1.10) -- (2.36,1.10);
\node[losslabflat,anchor=south,inner sep=1pt] at (2.10,1.17) {$\symResid$};
\MemColV{(2.40,1.34)}{4}{0.12}
\node[mathlab,anchor=north] at (2.46,0.82) {$\symMemNext$};
\node[minicap,anchor=west] at (2.60,1.10) {least squares,\\closed form};

% ---- 4. the fitted parameters -------------------------------------------
\draw[thinflow] (2.10,1.02) -- (2.10,0.80) -- (2.00,0.80) -- (2.00,0.68);
\MatBlock{(1.62,0.66)}{4}{4}{0.08}{clLatentD}
\node[mathlab,anchor=north] at (1.78,0.30) {$\symKhat$};
\MatBlock{(2.06,0.66)}{4}{1}{0.08}{ours}
\node[mathlab,anchor=north] at (2.10,0.30) {$\symBhat$};
\MatBlock{(2.36,0.66)}{4}{1}{0.08}{black}
\node[mathlab,anchor=north] at (2.40,0.30) {$\symbhat$};
 \end{scope}
\begin{scope}[shift={(4.985,0.00)}]     % modD_control.tex — module D: predictive test-time control.
% Draws eq:rollout (latent rollout per candidate, no LLM call), the
% closed-form expansion's per-step command injection (eq:closed_form_rollout),
% and eq:selection (cost + argmin over the candidate set).
% Canvas 4.735 x 3.03; top-left 2.2 x 0.30 left empty for the corner tag.
\path (0,0) rectangle (4.735,3.03);

% ---------------- 1. candidate set C_t --------------------------------
\node[mathlab,anchor=south] at (0.51,2.57) {$\symCmdSet$};
\ScalarCell{(0.12,2.55)}{ours}{}
\ScalarCell{(0.40,2.55)}{bad}{}
\ScalarCell{(0.68,2.55)}{good}{}

% three thin lines: the same chain is rolled out once per candidate
\draw[draw=ours,line width=0.4pt] (0.23,2.33) -- (0.14,2.20);
\draw[draw=bad,line width=0.4pt]  (0.51,2.33) -- (0.14,2.20);
\draw[draw=good,line width=0.4pt] (0.79,2.33) -- (0.14,2.20);

% ---------------- 2. latent rollout without an LLM call ---------------
% xihat_{t+h|t} = K xihat_{t+h-1|t} + B c + b   (eq:rollout, eq:closed_form_rollout)
% label the chain ONCE at each end; the fading colour marks successive h.
\LatCol{(0.05,2.20)}{3}{0.18}
\node[mathlab,anchor=north] at (0.14,1.62) {$\symLat$};
\LatColVF{(0.69,2.20)}{3}{0.18}{78}
\LatColVF{(1.33,2.20)}{3}{0.18}{58}
\LatColVF{(1.97,2.20)}{3}{0.18}{42}
\node[mathlab,anchor=north] at (2.06,1.62) {$\symLatRoll$};

\draw[thinflow] (0.25,1.93) -- (0.67,1.93);
\draw[thinflow] (0.89,1.93) -- (1.31,1.93);
\node[mathlab,anchor=south] at (1.10,1.99) {$\symK$};
\draw[thinflow] (1.53,1.93) -- (1.95,1.93);

% one instance of the per-step command injection is enough: the three
% arrows are identical, so a single Bc (ours) + b (grey) reads as "every
% step adds the command".
\MatBlock{(1.02,1.86)}{3}{1}{0.07}{ours}
\MatBlock{(1.10,1.86)}{3}{1}{0.07}{black}

\DimBraceB{(0.05,1.30)}{(2.15,1.30)}{$\symStep = 1,\ldots,\symHor$}
\node[minicap] at (1.10,0.85) {no LLM call};

% ---------------- 3. decode to an attribute (readout, once) -----------
% sits on the chain's own line, to the right of it -- not stacked below.
\draw[thinflow] (2.17,1.93) -- (2.36,1.83);
\begin{scope}[shift={(2.42,1.80)},scale=0.26]
  \DecNet{mdDec}{(0,0)}
\end{scope}
\draw[thinflow] (2.72,1.80) -- (2.75,1.78);
\RowSelector{(2.75,1.87)}{3}{0.06}
\node[mathlab,anchor=south] at (2.84,1.87) {$\symReadout$};
\draw[thinflow] (2.93,1.78) -- (2.95,1.55);

% ---------------- 4. predicted trajectory vs. target -------------------
\MiniAxes{(2.95,0.55)}{1.50}{1.10}
\node[mathlab,anchor=west] at (4.47,0.55) {$\symStep$};
\node[mathlab,anchor=south east] at (2.98,1.66) {$\symYroll$};
\draw[dashed,draw=good,line width=0.5pt] (2.95,1.10) -- (4.45,1.10);
\node[mathlab,anchor=west,text=good] at (4.47,1.10) {$\symTarget$};
\draw[draw=bad,line width=0.6pt]  (3.00,0.70) -- (3.45,1.40) -- (3.90,1.50) -- (4.40,1.35);
\draw[draw=good,line width=0.6pt] (3.00,0.70) -- (3.45,0.95) -- (3.90,1.05) -- (4.40,1.10);
\draw[draw=ours,line width=0.6pt] (3.00,0.70) -- (3.45,0.78) -- (3.90,0.85) -- (4.40,0.92);
\node[minicap] at (4.47,0.92) {$\symBound$};

% ---------------- 5. cost and the choice (above the plot) --------------
\node[mathlab,anchor=south] at (3.30,2.66) {$\symCost$};
\fill[ours] (2.95,2.48) rectangle (3.50,2.60);
\fill[bad]  (2.95,2.30) rectangle (3.80,2.42);
\fill[good] (2.95,2.12) rectangle (3.20,2.24);
\ScalarCell{(4.15,2.24)}{ours}{$\symCmdStar$}
\draw[oursflow] (4.15,2.13) -- (3.20,2.18);

% ---------------- 6. receding horizon (routed clear of the plot/bars) --
\draw[thinflow] (4.26,2.02) -- (4.70,2.02) -- (4.70,0.15) -- (0.06,0.15) -- (0.06,0.33);
\node[minicap,anchor=south] at (2.30,0.16) {execute, observe, replan};
  \end{scope}

% ---- cross-module relation (drawn): A's finite-memory state s_t is what
% B's dynamics model takes in. Both modules place their MemColV for s_t
% with its TOP at local y=2.09 (same column, same coordinate, in both
% frames), so a level arrow across the 0.25cm gap reads cleanly and does
% not cross either module's own content. Shifted up by \RowOneY (was the
% old fixed 3.28) since row 1 now sits higher. ----
% v3 (post CP2-redo): this is now the ONLY cross-module arrow in the right
% panel (G9). The old "(feeds D)" caption (B's yhat "feeds" D's rollout)
% is dropped, not replaced with a line -- the plan's rule for this version
% is that every other cross-module relation is carried by glyph identity
% alone (same shape, same colour, same symbol in both panels), never by a
% second drawn arrow or caption.
\draw[thinflow] (1.91,{2.09+\RowOneY}) -- (3.81,{2.09+\RowOneY});

% ---- C trains B (caption instead of a line; the two panels are not
% adjacent at this row spacing) ----
% DROPPED (step 3 rev 2): C's top strip has no genuinely free band for it --
% see FIG_STATE.md. C's own bottom caption "updates (phi,psi,K,B,b);
% theta frozen" already names B as one of the trained parameters.
